\section{背景知识：地址空间布局}

在 32 位 x86 系统中，虚拟地址空间共 4\,GiB（\hex{00000000}--\hex{FFFFFFFF}）。
High-Half Kernel 的经典划分方式是：

\begin{itemize}
  \item \textbf{低 3\,GiB}（\hex{00000000}--\hex{BFFFFFFF}）：用户态进程
  \item \textbf{高 1\,GiB}（\hex{C0000000}--\hex{FFFFFFFF}）：内核
\end{itemize}

\begin{figure}[H]
  \centering
  % figures/ch09/fig01-high-half-layout.tex — auto-extracted from chapter source
\begin{tikzpicture}[
  block/.style={draw, minimum width=6cm, font=\small, anchor=south west},
  addr/.style={font=\ttfamily\scriptsize\color{gray}, anchor=east},
]
  % 用户态区域
  \node[block, minimum height=4.5cm, fill=blue!8] (user) at (0, 0)
    {\parbox{5.5cm}{\centering\small\textbf{用户态空间}\\[2pt]
     \texttt{0x00000000 -- 0xBFFFFFFF}\\[2pt]
     进程代码、数据、堆、栈\\
     \textit{（3\,GiB，每进程独立）}}};

  % 内核区域
  \node[block, minimum height=1.5cm, fill=red!10] (kern) at (0, 4.5)
    {\parbox{5.5cm}{\centering\small\textbf{内核空间}\\[2pt]
     \texttt{0xC0000000 -- 0xFFFFFFFF}\\
     \textit{（1\,GiB，所有进程共享）}}};

  % 地址标注
  \node[addr] at (-0.2, 0)    {\hex{00000000}};
  \node[addr] at (-0.2, 4.5)  {\hex{C0000000}};
  \node[addr] at (-0.2, 6.0)  {\hex{FFFFFFFF}};

  % PD index 标注
  \draw[decorate, decoration={brace, amplitude=5pt, mirror}]
    (6.2, 0) -- (6.2, 4.5) node[midway, right=8pt, font=\scriptsize] {PD[0..767]};
  \draw[decorate, decoration={brace, amplitude=5pt, mirror}]
    (6.2, 4.5) -- (6.2, 6.0) node[midway, right=8pt, font=\scriptsize] {PD[768..1023]};

  % 内核子区域
  \draw[dashed, gray] (0.1, 5.25) -- (5.9, 5.25);
  \node[font=\scriptsize, anchor=west] at (0.3, 5.6) {动态分配区 / 堆};
  \node[font=\scriptsize, anchor=west] at (0.3, 4.9) {直接映射区};
\end{tikzpicture}

  \caption{High-Half Kernel 地址空间布局（32 位 x86）}
  \label{fig:high-half-layout}
\end{figure}

\begin{whybox}
为什么内核占高 1\,GiB 而不是低 1\,GiB？

将内核放在高地址有几个好处：
\begin{enumerate}
  \item 用户进程总是从 \hex{00000000} 附近加载（ELF 默认基址 \hex{08048000}），不用跳过内核区域
  \item 内核切换进程时只需替换低 3\,GiB 的页表（PD[0..767]），高 1\,GiB（PD[768..1023]）在所有进程中共享，切换代价小
  \item Linux 2.x/3.x、Windows NT 均采用此方案，是业界惯例
\end{enumerate}
\end{whybox}

% ============================================================================

